%% Question bank chapter file - generated by bank_tools.py build (do not edit while the bank is still being built:
%% it is regenerated from the per-test data in Question Bank\data). Every question is one bankq block:
%%   \begin{bankq}[id=..., source={...}, paper=.., qno=.., type=..., subtopic={...}, answer={...}, answer_src=key|solved,
%%                 working={...}, difficulty=1-10, diff_parts={K S I C F}, diff_note={...},
%%                 group=<passage id or empty>, also_in={duplicate copies in other tests}]
%% difficulty = min(10, K+S+I+C+F-1): K concepts 1-3, S steps 1-4, I insight 0-2, C computation 0-2, F format 0-1
%%   ...question LaTeX, options, figures (\bfig / \bfigi -> img/)...
%%   \end{bankq}
%% Passages shared by several questions are bankpassage blocks; their questions carry group=<passage id>.
\documentclass[10pt]{article}
\usepackage[a4paper,margin=1.8cm]{geometry}
\usepackage[T1]{fontenc}
\usepackage{amsmath,amssymb,mathtools}
\usepackage{graphicx}
\usepackage[export]{adjustbox}
\usepackage{xcolor,needspace,keyval}
\usepackage{enumitem}
\usepackage{booktabs,array,tabularx}
\usepackage[hidelinks]{hyperref}
\setlength{\parindent}{0pt}
\setlength{\parskip}{3pt}
\makeatletter
\define@key{bq}{id}{\def\bq@id{#1}}\define@key{bq}{source}{\def\bq@src{#1}}\define@key{bq}{paper}{\def\bq@paper{#1}}
\define@key{bq}{qno}{\def\bq@qno{#1}}\define@key{bq}{type}{\def\bq@type{#1}}\define@key{bq}{subtopic}{\def\bq@sub{#1}}
\define@key{bq}{answer}{\def\bq@ans{#1}}\define@key{bq}{answer_src}{\def\bq@asrc{#1}}\define@key{bq}{working}{\def\bq@work{#1}}
\define@key{bq}{group}{\def\bq@group{#1}}\define@key{bq}{also_in}{\def\bq@also{#1}}
\define@key{bq}{difficulty}{\def\bq@diff{#1}}\define@key{bq}{diff_parts}{\def\bq@dparts{#1}}\define@key{bq}{diff_note}{\def\bq@dnote{#1}}
\newenvironment{bankq}[1][]{\def\bq@work{}\def\bq@also{}\def\bq@group{}\def\bq@diff{?}\def\bq@dparts{}\def\bq@dnote{}\setkeys{bq}{#1}%
  \par\needspace{4\baselineskip}\medskip\noindent{\small\textbf{\bq@id}\hfill\color{gray}[\bq@src, Paper \bq@paper, Q\bq@qno]
  \bq@type\ $\cdot$ \bq@sub\ $\cdot$ \textbf{D\bq@diff}}\par\nopagebreak\relax}
  {\par{\small\color{teal}Answer: \textbf{\bq@ans}\ifx\bq@work\@empty\else\ (solved: \bq@work)\fi
  \ifx\bq@also\@empty\else\par\color{gray}Also in: \bq@also\fi}\par}
\define@key{bp}{id}{\def\bp@id{#1}}\define@key{bp}{source}{\def\bp@src{#1}}\define@key{bp}{questions}{\def\bp@qs{#1}}
\newenvironment{bankpassage}[1][]{\setkeys{bp}{#1}\par\needspace{4\baselineskip}\bigskip\noindent
  {\small\textbf{Passage \bp@id}\hfill\color{gray}[\bp@src; questions \bp@qs]}\par\nopagebreak\relax\itshape}{\par}
\makeatother
\newcommand{\bfig}[2]{\par\smallskip{\centering\includegraphics[width=#1pt,max width=0.62\linewidth,max height=0.32\textheight]{img/#2}\par}\smallskip}
\newcommand{\bfigi}[2]{\includegraphics[width=#1pt,max width=\linewidth,max height=0.32\textheight,valign=c]{img/#2}}
\newcommand{\bmissing}[1]{\textit{[figure missing in the source conversion: \detokenize{#1}]}}
\newenvironment{qoptions}{\begin{itemize}[leftmargin=1.8em,itemsep=1pt,topsep=2pt,label={}]}{\end{itemize}}
\newcommand{\opt}[2]{\item[(#1)] #2}
\providecommand{\fwstart}{}\providecommand{\fwend}{}\providecommand{\qfullsplit}{}
\newcolumntype{L}[1]{>{\raggedright\arraybackslash}p{#1}}
\newcolumntype{C}[1]{>{\centering\arraybackslash}p{#1}}
\newcolumntype{R}[1]{>{\raggedleft\arraybackslash}p{#1}}
\providecommand{\sectionhdr}[1]{}\providecommand{\subsectionhdr}[1]{}


\begin{document}
{\Large\bfseries Trig Ratios \& Transformation --- Question Bank}\par\smallskip
{\small Mathematics chapter M11. Questions from all converted test papers, grouped by source series. Answers: printed key where it matches the paper, otherwise solved (answer\_src=solved, with working). Exact repeats across papers appear once, with also\_in listing the other copies.}\par\medskip\hrule

\section*{TOAS papers}
\begin{bankpassage}[id=TOAS01-P1-MATH-S3-G3-4, source={TOAS - Test 01 (30 Jan 2022)}, questions={3-4}]
Let `$A$' is the number of solution(s) of the equation
\[x^{\log x\,\log 4}\times4^{\log x^3}+\dfrac{64}{4^{\log^2x}\times x^{3\log4}}=65,\]
$B$ is the value of $\sec^2x$ if $\dfrac{\sin^2x-2\cos^2x+1}{\sin^2x+2\cos^2x-1}=4$ and $C$ is the value of the expression $\cos2x-\dfrac{\sin x\cos3x}{\sin2x}$, (wherever defined).
\end{bankpassage}
\begin{bankq}[id=TOAS01-P1-MATH-S3-Q03, source={TOAS - Test 01 (30 Jan 2022)}, paper=1, qno=3, type=numerical, subtopic={Identities}, answer={4}, answer_src=key, difficulty=4, diff_parts={K1 S2 I1 C0 F1}, diff_note={Rewrite in cos^2x to get sec^2=7/2; triple-angle simplification gives 1/2}, group=TOAS01-P1-MATH-S3-G3-4]
The value of $(B+C)$ is -
\end{bankq}
\begin{bankq}[id=TOAS01-P1-MATH-S3-Q04, source={TOAS - Test 01 (30 Jan 2022)}, paper=1, qno=4, type=numerical, subtopic={Identities}, answer={0}, answer_src=key, difficulty=9, diff_parts={K3 S4 I1 C1 F1}, diff_note={Log equation to u+64/u=65 gives A=4 roots, then least integral value of the quadratic with B, C}, group=TOAS01-P1-MATH-S3-G3-4]
Least integral value of $(Ax^2+Bx+C)$ is -
\end{bankq}
\begin{bankq}[id=TOAS03-P1-MATH-S1-Q04, source={TOAS - Test 03 (27 Feb 2022)}, paper=1, qno=4, type=single, subtopic={Compound and multiple angles}, answer={A}, answer_src=key, difficulty=4, diff_parts={K1 S3 I1 C0 F0}, diff_note={Reduce to (t+1)^2+2 with t=sin x - cos x in [-sqrt2, sqrt2]}]
The greatest value of the function $f(x)=(\sin x-\cos x)^2+(\sin x+1)^2+(\cos x-1)^2$ is :-
\begin{qoptions}
\opt{A}{$5+2\sqrt2$}
\opt{B}{$5-2\sqrt2$}
\opt{C}{24}
\opt{D}{16}
\end{qoptions}
\end{bankq}
\begin{bankpassage}[id=TOAS04-P2-MATH-S2-G7-8, source={TOAS - Test 04 (15 Mar 2022)}, questions={7-8}]
If $\tan^2\alpha\tan^2\beta+\tan^2\beta\tan^2\gamma+\tan^2\gamma\tan^2\alpha+2\tan^2\alpha\tan^2\beta\tan^2\gamma=1$.\\
(where $\alpha,\beta,\gamma\in R-\left\{\dfrac{n\pi}2\right\}$, $n\in I$)
\end{bankpassage}
\begin{bankq}[id=TOAS04-P2-MATH-S2-Q07, source={TOAS - Test 04 (15 Mar 2022)}, paper=2, qno=7, type=single, subtopic={Identities}, answer={D}, answer_src=key, difficulty=6, diff_parts={K1 S3 I2 C1 F0}, diff_note={Given relation in tan^2 gives product minus sum of cot^2 equal to 2}, group=TOAS04-P2-MATH-S2-G7-8]
Value of $\left(\Pi\cot^2\alpha-\Sigma\cot^2\alpha\right)$ is equal to -
\begin{qoptions}
\opt{A}{0}
\opt{B}{$-1$}
\opt{C}{1}
\opt{D}{2}
\end{qoptions}
\end{bankq}
\begin{bankq}[id=TOAS04-P2-MATH-S2-Q08, source={TOAS - Test 04 (15 Mar 2022)}, paper=2, qno=8, type=single, subtopic={Identities}, answer={C}, answer_src=key, difficulty=6, diff_parts={K2 S3 I2 C0 F0}, diff_note={Same relation yields cos 2a + cos 2b + cos 2c = 1}, group=TOAS04-P2-MATH-S2-G7-8]
Value of $\cos2\alpha+\cos2\beta+\cos2\gamma$ is equal to -
\begin{qoptions}
\opt{A}{$-1$}
\opt{B}{0}
\opt{C}{1}
\opt{D}{2}
\end{qoptions}
\end{bankq}
\begin{bankpassage}[id=TOAS04-P2-MATH-S3-G3-4, source={TOAS - Test 04 (15 Mar 2022)}, questions={3-4}]
Let $a=\sin10^\circ$, $b=\sin50^\circ$, $c=\sin70^\circ$.
\end{bankpassage}
\begin{bankq}[id=TOAS04-P2-MATH-S3-Q03, source={TOAS - Test 04 (15 Mar 2022)}, paper=2, qno=3, type=numerical, subtopic={Transformations}, answer={0.75}, answer_src=key, difficulty=8, diff_parts={K2 S3 I2 C1 F1}, diff_note={a+b=c by sum to product; abc=1/8 and c^2-ab=3/4 give the value}, group=TOAS04-P2-MATH-S3-G3-4]
Value of \[abc\left(\dfrac{a+b}{c}\right)\left(\dfrac1a+\dfrac1b-\dfrac1c\right)\] is equal to -
\end{bankq}
\begin{bankq}[id=TOAS04-P2-MATH-S3-Q04, source={TOAS - Test 04 (15 Mar 2022)}, paper=2, qno=4, type=numerical, subtopic={Transformations}, answer={2.00}, answer_src=key, difficulty=6, diff_parts={K1 S3 I2 C0 F1}, diff_note={With a+b-c=0 the identity x^4+y^4+z^4 = 2 e2^2 gives ratio 2}, group=TOAS04-P2-MATH-S3-G3-4]
Value of $\dfrac{a^4+b^4+c^4}{(ab)^2+(bc)^2+(ca)^2}$ is -
\end{bankq}
\begin{bankq}[id=TOAS05-P2-MATH-S1-Q14, source={TOAS - Test 05 (03 Apr 2022)}, paper=2, qno=14, type=single, subtopic={Compound and multiple angles}, answer={C}, answer_src=key, difficulty=2, diff_parts={K1 S2 I0 C0 F0}, diff_note={Amplitude sqrt(k^2+1)=5 so k=+-sqrt24, count integers between}]
Range of $f(x)=k\sin x+\cos x$ is $[-5,5]$, then total number of integers between the two possible value of `$k$' is-
\begin{qoptions}
\opt{A}{11}
\opt{B}{10}
\opt{C}{9}
\opt{D}{5}
\end{qoptions}
\end{bankq}
\begin{bankq}[id=TOAS05-P2-MATH-S1-Q15, source={TOAS - Test 05 (03 Apr 2022)}, paper=2, qno=15, type=single, subtopic={Identities}, answer={C}, answer_src=key, difficulty=3, diff_parts={K1 S2 I1 C0 F0}, diff_note={Square and add the two relations}]
If $\sin A+\cos B=a$ and $\sin B+\cos A=b$, then $\sin(A+B)$ is equal to-
\begin{qoptions}
\opt{A}{$\dfrac{a^2+b^2}2$}
\opt{B}{$\dfrac{a^2-b^2+2}2$}
\opt{C}{$\dfrac{a^2+b^2-2}2$}
\opt{D}{$\dfrac{a^2+b^2+2}2$}
\end{qoptions}
\end{bankq}
\begin{bankq}[id=TOAS05-P2-MATH-S1-Q16, source={TOAS - Test 05 (03 Apr 2022)}, paper=2, qno=16, type=single, subtopic={Compound and multiple angles}, answer={D}, answer_src=key, difficulty=3, diff_parts={K1 S2 I1 C0 F0}, diff_note={Multiply by sin(pi/5) and use the doubling identity}]
The value of $\cos\dfrac\pi5\cos\dfrac{2\pi}5\cos\dfrac{4\pi}5\cos\dfrac{8\pi}5$ is equal to-
\begin{qoptions}
\opt{A}{$\dfrac1{16}$}
\opt{B}{$\dfrac14$}
\opt{C}{$-\dfrac18$}
\opt{D}{$-\dfrac1{16}$}
\end{qoptions}
\end{bankq}
\end{document}
